\section{Описание}
Требуется выбрать самый дешёвый набор пищевых добавок, по измерениям воздействия которых можно однозначно найти неизвестные коэффициенты $c_1, \dots, c_N$, либо сообщить, что это невозможно.

Каждая добавка описывается вектором $a = (a_1, \dots, a_N)$ и ценой. Измерив воздействие добавки $i$, мы получаем линейное уравнение $c_1 a_{i1} + \dots + c_N a_{iN} = y_i$. Для выбранного набора добавок получается система линейных уравнений относительно $c_1, \dots, c_N$, и она имеет единственное решение тогда и только тогда, когда ранг матрицы из векторов выбранных добавок равен $N$. Значит, коэффициенты можно определить, только если ранг матрицы всех $M$ добавок равен $N$, иначе нужно вывести $-1$.

Если ранг равен $N$, то среди добавок нужно найти набор минимальной стоимости, ранг которого равен $N$. Цены неотрицательны, поэтому из любого такого набора можно выбросить линейно зависимые добавки, не увеличив стоимость. Следовательно, оптимальный набор состоит из ровно $N$ линейно независимых векторов, то есть является базисом.

Множество строк матрицы вместе с семейством линейно независимых подмножеств образует матроид, а для задачи о базисе минимального веса в матроиде жадный алгоритм даёт оптимальный ответ \cite{Kormen}. Интуитивно это следует из того, что добавка, лежащая в линейной оболочке более дешёвых уже выбранных добавок, ничего не добавляет к рангу, и её можно заменить более дешёвыми. Поэтому алгоритм такой:
\begin{enumerate}
    \item отсортировать добавки по возрастанию цены;
    \item перебирать добавки в этом порядке и брать очередную, если она линейно независима от уже взятых;
    \item остановиться, когда набрано $N$ добавок. Если добавки закончились раньше, вывести $-1$.
\end{enumerate}

Линейная независимость проверяется методом Гаусса. Для каждого столбца $k$ хранится индекс строки \texttt{basis[k]}, у которой первый ненулевой элемент стоит в столбце $k$ и равен единице. Очередная строка последовательно приводится по этим базисным строкам: для каждого столбца $k$, в котором элемент строки не равен нулю, либо вычитается базисная строка (если она есть), либо строка сама становится базисной строкой для столбца $k$ после деления на этот элемент. Если строка свелась к нулю, она линейно зависима от выбранных и пропускается. Столбец цены в вычитании не участвует.

Так как цены~--- целые числа от 0 до 50, сортировка выполняется сортировкой подсчётом. Она устойчива и запоминает исходные номера добавок, которые нужны для вывода. Найденные номера затем тоже сортируются подсчётом по возрастанию.

Сложность по времени: сортировка $O(M + C)$, где $C = 50$~--- число возможных цен, приведение одной строки $O(N^2)$ (не более $N$ вычитаний по $O(N)$ каждое), всего $O(MN^2)$, а так как $M \le N$, то не более $O(N^3)$. Сложность по памяти~--- $O(MN)$.

Вычисления ведутся в числах с плавающей точкой (\texttt{double}), ноль определяется сравнением с $\varepsilon = 10^{-10}$. Для небольших $N$ и коэффициентов не больше 50 этого достаточно, но при больших $N$ накопление ошибок округления может привести к неверному определению ранга. Точные варианты~--- вычисления в дробях или по модулю большого простого числа.

\pagebreak

\section{Исходный код}
Исходные данные хранятся в матрице \texttt{input} размера $M \times (N+1)$: последний элемент каждой строки~--- цена. Функция \texttt{counting\_sort\_input} устойчиво сортирует строки по цене сортировкой подсчётом и возвращает массив \texttt{permutations}, где \texttt{permutations[j]}~--- исходный номер строки, оказавшейся на позиции $j$. Исходная матрица после сортировки освобождается.

Затем строки обрабатываются в порядке возрастания цены. Массив \texttt{basis} размера $N$ хранит для каждого столбца номер базисной строки (или $-1$). Для очередной строки внутренний цикл идёт по столбцам: нулевые элементы пропускаются, при отсутствии базисной строки в столбце текущая строка становится базисной (нормируется делением на ведущий элемент), иначе из неё вычитается базисная строка. Внутренние циклы по $j$ идут от $N-1$ к $k$, чтобы элемент $k$, на который умножаются и делятся остальные, изменялся последним. Переменная \texttt{full\_basis} позволяет остановиться, когда базис уже полон.

После обработки собираются исходные номера базисных строк. Если их меньше $N$, выводится $-1$, иначе номера сортируются функцией \texttt{counting\_sort\_ans} и выводятся.

\begin{lstlisting}[language=C++]
#include <algorithm>
#include <iostream>
#include <utility>
#include <vector>

#define MAX_PRICE 50
#define EPS 1e-10

std::vector<int> counting_sort_input(std::vector<std::vector<double>>& v, std::vector<std::vector<double>>& res) {
    std::vector<int> count(MAX_PRICE + 1);
    for (int i = 0; i < v.size(); ++i) {
        ++count[v[i].back()];
    }
    for (int i = 1; i < count.size(); ++i) {
        count[i] += count[i - 1];
    }
    std::vector<int> permutations(v.size());
    res.resize(v.size());
    for (int i = v.size() - 1; i >= 0; --i) {
        int j = --count[v[i].back()];
        res[j] = std::move(v[i]);
        permutations[j] = i;
    }
    return permutations;
}

std::vector<int> counting_sort_ans(std::vector<int>& v) {
    std::vector<int> count(1 + *std::max_element(v.begin(), v.end()));
    for (int i = 0; i < v.size(); ++i) {
        ++count[v[i]];
    }
    for (int i = 1; i < count.size(); ++i) {
        count[i] += count[i - 1];
    }
    std::vector<int> res(v.size());
    for (int i = v.size() - 1; i >= 0; --i) {
        int j = --count[v[i]];
        res[j] = v[i];
    }
    return res;
}

int main(void) {
    int m, n;
    std::cin >> m >> n;

    std::vector<std::vector<double>> input;
    input.resize(m);
    for (int i = 0; i < m; ++i) {
        input[i].resize(n + 1);
        for (int j = 0; j <= n; ++j) {
            std::cin >> input[i][j];
        }
    }

    std::vector<std::vector<double>> bads;
    auto permutations = counting_sort_input(input, bads);
    input.clear();

    std::vector<int> basis; // [индекс базисного элемента] = индекс строки
    basis.resize(n, -1);

    for (int i = 0; i < m; ++i) {
        bool full_basis = true;
        // вычитаем из новой строки все что есть в базисе
        for (int k = 0; k < n; ++k) {
            full_basis &= (basis[k] != -1);

            if (std::abs(bads[i][k]) < EPS) {
                continue;
            }

            if (basis[k] == -1) {
                basis[k] = i;
                for (int j = n-1; j >= k; --j) {
                    bads[i][j] /= bads[i][k];
                }
                break;
            }

            for (int j = n-1; j >= k; --j) {
                bads[i][j] -= bads[i][k] * bads[basis[k]][j];
            }
        }

        if (full_basis) {
            break;
        }
    }

    std::vector<int> ans;
    for (int i = 0; i < n; ++i) {
        if (basis[i] == -1) {
            break;
        }
        ans.push_back(permutations[basis[i]] + 1);
    }

    if (ans.size() < n) {
        std::cout << -1;
    } else {
        ans = counting_sort_ans(ans);
        for (auto& i : ans) {
            std::cout << i << " ";
        }
    }
}
\end{lstlisting}

\section{Консоль}
\begin{verbatim}
laba-8> cmake -S . -B build
laba-8> cmake --build build --config Release
laba-8> type t1
3 3
1 0 0 5
0 1 0 3
0 0 1 4
laba-8> build\Release\main.exe < t1
1 2 3
laba-8> type t2
3 2
1 1 2
1 0 5
0 1 1
laba-8> build\Release\main.exe < t2
1 3
laba-8> type t3
2 2
1 1 3
2 2 4
laba-8> build\Release\main.exe < t3
-1
\end{verbatim}

Пояснения к тестам. В первом тесте все три вектора линейно независимы, поэтому нужны все добавки. Во втором тесте сначала берётся самая дешёвая добавка 3 (цена 1), затем добавка 1 (цена 2), которая независима от неё; дорогая добавка 2 (цена 5) уже не нужна, так как ранг равен 2. В третьем тесте добавка 2 пропорциональна добавке 1, ранг матрицы равен 1 при $N = 2$, поэтому коэффициенты определить нельзя.
\pagebreak